%% =============================================================================
%% PAPER 1 - Section 7: Conclusion
%% =============================================================================
We document positive and economically significant alpha in three
European fixed-income arbitrage strategies---the BTP~Italia
inflation-linked basis, the corporate CDS--bond basis, and the
iTraxx index skew---and show that this alpha is not compensation
for the identified systematic risk exposures, whether the controls
are entered as classical benchmark factors or as data-driven
selections from a large candidate universe, and it survives when each
benchmark hedge is re-estimated out of sample from past-only loadings. The mispricings are
linked by a common factor that loads on intermediary balance-sheet
variables (and not on macroeconomic fundamentals), and their
persistence at the level of each series reflects the
funding-dependence of the arbitrage trade: the CDS--Bond basis,
the only strategy requiring continuous repo financing of the cash
bond, is the most persistent, while the BTP~Italia and the iTraxx
skew mean-revert more rapidly and are essentially identical to each
other. The three correlations exhibit an asymmetric
pattern that mirrors the structural divide between institutional
and retail holder bases: positive and rising in stress within the
institutional pair (CDS--Bond and iTraxx), and negative across
holder bases (BTP~Italia vs.\ either institutional strategy). The
slow-moving capital literature
\citep{duffie2010presidential, brunnermeier2009market,
mitchell2012arbitrage}, combined with the post-crisis regulatory
environment documented by \citet{boyarchenko2016trends}, provides
an economic explanation for why these mispricings persist.

The findings have implications for two strands of the literature.
For the limits-to-arbitrage literature, the European setting
provides a natural test of the slow-moving capital mechanism in a
market with a structurally heterogeneous holder base, and the
investor-base-segmentation pattern in cross-strategy correlations
offers independent evidence that the geography of arbitrage
capital---rather than asset-class labels---determines cross-market
co-movement. For the fixed-income asset pricing literature, supervised selection
on a large candidate universe recovers strategy-specific risk drivers
that the standard benchmark factor sets miss, yet economically
significant alpha persists for the BTP~Italia and CDS--Bond
strategies---indicating that the residual return is genuinely
unspanned rather than compensation for an overlooked systematic
factor.